%% ============================================================================
%%  Response to Reviewers — IEEE Access, MS #Access-2026-27603
%%
%%  Manuscript: Stochastic ODE-Guard — A Neural SDE Framework with
%%              Anti-Concentration Bounds for Adversarially Robust Network
%%              Intrusion Detection.
%%
%%  Compile with pdflatex (self-contained; no external .bib required).
%%  Upload as "Author's Response Files" on the IEEE Author Portal.
%% ============================================================================
\documentclass[11pt,letterpaper]{article}

\usepackage[margin=1in]{geometry}
\usepackage[T1]{fontenc}
\usepackage{lmodern}
\usepackage{amsmath,amssymb,amsfonts,mathtools}
\usepackage{booktabs}
\usepackage{array}
\usepackage{longtable}
\usepackage{tabularx}
\usepackage{enumitem}
\usepackage{xcolor}
\usepackage[colorlinks=true,linkcolor=blue!60!black,urlcolor=blue!60!black,citecolor=blue!60!black]{hyperref}
\usepackage{fancyhdr}
\usepackage{titlesec}

%% ----- Layout -----
\setlength{\parindent}{0pt}
\setlength{\parskip}{6pt}
\pagestyle{fancy}
\fancyhf{}
\lhead{\small IEEE Access — MS \#Access-2026-27603}
\rhead{\small Response to Reviewers}
\cfoot{\thepage}

\titleformat{\section}{\normalfont\Large\bfseries\color{blue!60!black}}{\thesection}{0.7em}{}
\titleformat{\subsection}{\normalfont\large\bfseries}{\thesubsection}{0.7em}{}

%% ----- Colour-coded blocks -----
\definecolor{ReviewerColor}{RGB}{40,40,40}
\definecolor{QuoteBg}{RGB}{245,245,235}
\definecolor{ResponseAccent}{RGB}{10,90,150}

\newcommand{\rev}[2]{%
  \subsection*{\normalsize\textbf{\color{ResponseAccent}#1}\quad\textmd{\textit{#2}}}%
}

%% Simple quote box used throughout:
\newcommand{\rquote}[1]{%
  \begin{center}%
    \colorbox{QuoteBg}{\parbox{0.94\linewidth}{\vspace{2pt}\itshape\small\textbf{\upshape Reviewer:}\ #1\vspace{2pt}}}%
  \end{center}%
}

\newcommand{\resp}{\par\textbf{\color{ResponseAccent}Response.}\ }
\newcommand{\action}{\par\textbf{\color{ResponseAccent}Action.}\ }

%% ============================================================================
\title{\vspace{-1.4cm}
       \textbf{Response to Reviewers}\\[4pt]
       \normalsize Manuscript ID: \texttt{Access-2026-27603}\\
       \normalsize \emph{Stochastic ODE-Guard: A Neural Stochastic
       Differential Equation Framework with Anti-Concentration Bounds for
       Adversarially Robust Network Intrusion Detection}}
\author{R.\ N.\ Anaedevha, A.\ G.\ Trofimov, and Y.\ V.\ Borodachev\\
        \normalsize National Research Nuclear University MEPhI, Moscow, Russia\\
        \normalsize \texttt{ar006@campus.mephi.ru}}
\date{\normalsize Resubmission version~4 --- \today}
%% ============================================================================

\begin{document}
\maketitle
\thispagestyle{fancy}
\vspace{-0.6cm}

\section*{Preamble}
We thank both reviewers for their careful and constructive reading. Every
concern raised in the decision letter has been addressed in the revised
manuscript (highlighted in the \emph{Highlighted PDF} submission) and,
where appropriate, by new implementation modules in the reproducibility
repository. In this document we quote each concern, describe the
response, and cite the section, table, or code file where the change
lives. The implementation additions ship in the same commit as the
manuscript revision and are runnable with
\texttt{bash scripts/reproduce\_revision.sh}. The reproducibility
artefact is publicly hosted at
\url{https://github.com/rogerpanel/CV/tree/main/SODE-Guard}.

Table~\ref{tab:overview} indexes each reviewer point to the location
of the fix in the revised manuscript.

\begin{table}[h]
\centering\small
\caption{Cross-reference: reviewer concern $\rightarrow$ location of fix.}
\label{tab:overview}
\begin{tabularx}{\linewidth}{@{}l l X@{}}
\toprule
\textbf{Point} & \textbf{Fix location} & \textbf{Nature of change} \\
\midrule
R1.1 & Lemma~3, Table~6, \texttt{src/theory/lipschitz.py} & Closed-form $L_g$ + Hoeffding empirical certificate \\
R1.2 & \S6.7, Table~10, \texttt{src/evaluation/reliability/commercial\_context.py} & Commercial baselines moved to a separate reference category \\
R1.3 & \S6.5.1, Table~5, \texttt{src/evaluation/reliability/smoothing\_sweep.py} & $\sigma$-swept randomised-smoothing baseline \\
R1.4 & Eq.~(7), Alg.~2, \texttt{src/theory/chaos\_degree.py} & Precise $\mathcal L_{\mathrm{AC}}$ + Hermite regression \\
R1.5 & Thm.~2 / Eq.~(2), \texttt{src/theory/pseudoinverse.py} & Moore--Penrose pseudo-inverse stated in-line \\
R1.6 & \S5.4, Table~7, \texttt{src/theory/pac\_bayes.py} & PAC-Bayes bound tied to ECE \\
\midrule
R2.1 & \S1 (contributions) & Novelty framed as principled composition \\
R2.2 & Prop.~5, App.~B, \texttt{src/theory/carbery\_wright.py} & Direction-corrected Carbery--Wright with lower CB \\
R2.3 & Def.~4, Cor.~1, App.~C & ``Certified radius'' $\rightarrow$ probabilistic robustness radius; argmax bound \\
R2.4 & Alg.~3, \texttt{src/theory/chaos\_degree.py} & Adaptive per-example chaos degree $d^\star_{p99}$ \\
R2.5 & Rem.~8 & Explicit contrast with randomised smoothing \\
R2.6--R2.9 & Assumption~1, App.~D--E, \texttt{src/theory/pseudoinverse.py}, \texttt{bel\_estimator.py} & Assumption audit + BEL bias / variance vs FD \\
R2.10--R2.11 & \S4.1, Table~9, \texttt{src/attacks/semantic/} & SDE on embedded state; problem-space feasibility projector \\
R2.12 & App.~F, \texttt{src/data/splits\_extra/dedup.py} & Dataset provenance + dedup + leakage report \\
R2.13 & \S6.6, Table~7, \texttt{src/data/splits\_extra/} & Temporal, host-disjoint, scenario-disjoint splits \\
R2.14 & Table~11 & Un-harmonised IIS3D per-constituent results \\
\bottomrule
\end{tabularx}
\end{table}

%% ============================================================================
\section{Reviewer 1}
%% ============================================================================

%% ---------------------------------------------------------------------------
\rev{R1.1}{$L_g$ assumption without proof / verification}
\rquote{Proposition 5 rests on the assumption that the truncated chaos
expansion satisfies an $L^2$ Lipschitz bound in $\delta$ (i.e.,
$\mathbb E[(g(x)-g(x+\delta))^2]^{1/2}\le L_g\|\delta\|$). This
assumption is stated without a constructive proof or empirical
verification beyond the ablation.}

\resp
The revised Proposition~5 no longer treats $L_g$ as an unknown
hyper-parameter. Two independent changes address the concern:
\begin{enumerate}[leftmargin=1.3em]
  \item \textbf{Constructive derivation.} Appendix~A now proves that,
        under the linear-growth and Lipschitz conditions of Theorem~1
        and the ellipticity floor $\lambda_0>0$, the Grönwall
        estimate on the first-variation process combined with the
        Itô isometry yields the closed-form bound
        \[
          L_g \le \frac{\|\psi_\eta\|_{\Lip}}{T}
                  \exp\!\bigl( (K_f + \tfrac12 K_g^2)\,T\bigr).
        \]
        This is now stated as \textbf{Lemma~3}.
  \item \textbf{Empirical certification.} Module
        \texttt{src/theory/lipschitz.py} draws $K=8$ random directions
        per test point, integrates $N=32$ Brownian paths, and returns
        a Hoeffding upper-confidence bound on $L_g$ at the 95\% level.
        The certified $\widehat L_g$ values now appear in
        \textbf{Table~6} of the revised manuscript, and Proposition~5
        uses the corresponding \emph{lower} confidence bound
        $L_{\mathrm{lo}}$ in the denominator.
\end{enumerate}
\action
Added Lemma~3, Table~6, and the new module cited above; every use of
$L_g$ in the manuscript is now anchored to the empirically-certified
$L_{\mathrm{lo}}$.

%% ---------------------------------------------------------------------------
\rev{R1.2}{Fair comparison with commercial baselines}
\rquote{The comparison with commercial signature-based engines (Snort,
Suricata) and the LLM-based filter (Llama Guard) is not entirely fair,
as these systems are not designed for the same threat model \dots the
Pareto frontier plot (Fig.~2) mixes fundamentally different paradigms.}

\resp
The revised paper separates the comparison into two categories
mirroring the taxonomy exposed by
\texttt{src/evaluation/reliability/commercial\_context.py}:
\begin{itemize}[leftmargin=1.3em]
  \item \emph{Threat-model-matched neural baselines} (E-GraphSAGE,
        RTIDS, CNN-LSTM, IDS-GraphMamba, SurrogateIDS-7B, SDE-TGNN)
        share the $\ell_\infty$ gradient-based threat model with
        SODE-Guard and drive Tables~1--4 and Figures~2--3.
  \item \emph{Reference (non-neural) baselines} (Snort~3~+~SnortML,
        Suricata~7, Llama~Guard~3) are moved to
        \textbf{Table~10} with an explicit ``different threat model''
        caption. All Pareto claims in the main text now refer only to
        the threat-model-matched set.
\end{itemize}
\action
Rewrote \S6.2, moved the reference baselines to \S6.7, and re-drew
the radar figure using only the matched set.

%% ---------------------------------------------------------------------------
\rev{R1.3}{Randomised-smoothing $\sigma$ tuning}
\rquote{The comparison against randomised smoothing uses a fixed
$\sigma=0.25$ without tuning; a sensitivity analysis over $\sigma$ would
strengthen the claim that the proposed certificate is indeed superior in
practice.}

\resp
Section~6.5.1 now sweeps $\sigma\in\{0.10, 0.15, 0.20, 0.25, 0.30,
0.40, 0.50\}$ for the Cohen \emph{et al.} baseline and reports the
\emph{best} $\sigma$ per benchmark. The sweep is implemented in
\texttt{src/evaluation/reliability/smoothing\_sweep.py} and its
numerical output is shipped in
\texttt{experiments/results/smoothing\_sensitivity.csv}. SODE-Guard
remains above the $\sigma$-tuned smoothing baseline at every
operational target radius $r\ge 0.02$; Figure~4b of the revised paper
plots the two curves against one another.
\action
Added Table~5 and Figure~4b.

%% ---------------------------------------------------------------------------
\rev{R1.4}{Anti-concentration regulariser: precise formula}
\rquote{The definition of $\mathcal L_{\mathrm{AC}}$ (Remark~7) is given
only informally; a precise formula with the Hermite-regression
procedure would help reproducibility.}

\resp
Equation~(7) of the revised paper gives the closed-form
regulariser
\[
  \mathcal L_{\mathrm{AC}}(\theta) = \frac{1}{|\mathcal B|}\!\!\sum_{\beta\in\mathcal B}\!
    \log\!\Bigl(1 + C\,d^\star(\beta / \hat s)^{1/d^\star}
        \cdot \tfrac{1}{B}\!\sum_i\!\sigma_\kappa(\beta-|m_i|)\Bigr),
\]
with $\sigma_\kappa$ the temperature-$\kappa=50$ sigmoid surrogate for
the indicator, $m_i$ the smoothed margin (top-1 minus top-2) on $P$
paths, and $\mathcal B=\{0.01,0.025,0.05,0.10\}$. The full
step-by-step Hermite-regression procedure is Algorithm~2, and the
underlying code is a single module,
\texttt{src/theory/chaos\_degree.py}, shared between training-time
regularisation and inference-time certification.
\action
Added Equation~(7), Algorithm~2, and the module cited above.

%% ---------------------------------------------------------------------------
\rev{R1.5}{$g^{-1}$ notation for non-square $g$}
\rquote{The notation in the BEL identity uses $g^{-1}$ even though $g$
is not square \dots this should be stated explicitly in the theorem
statement.}

\resp
Theorem~2 of the revised paper writes $g^{+}$ (Moore--Penrose
pseudo-inverse) throughout and defines it in-line as
\[
  g^{+}(x)=\bigl(g(x)^{\!\top}g(x)+\lambda_0 I_m\bigr)^{-1}g(x)^{\!\top}.
\]
The regularised form is well-conditioned by the ellipticity floor
$\lambda_0=10^{-3}$ (median condition number $8.9$; see Appendix~D).
The pseudo-inverse routine used by the BEL estimator and by
\texttt{src/theory/pseudoinverse.py} matches the mathematics
one-for-one.
\action
Replaced $g^{-1}$ with $g^{+}$ globally and added the definition in
Equation~(2).

%% ---------------------------------------------------------------------------
\rev{R1.6}{PAC-Bayes disconnected from results}
\rquote{The PAC-Bayes certificate is mentioned but not used in the
main results (e.g., Table~1); I suggest either connecting it more
tightly to the reported calibration errors or moving it to the
discussion section to avoid overloading the reader.}

\resp
Section~5.4 of the revised paper reports the Maurer PAC-Bayes-kl
population-risk certificate on all three benchmarks in \textbf{Table~7},
together with the empirical calibration error (ECE). The certificate
is computed by \texttt{src/theory/pac\_bayes.py} using the Reeb--Seldin
kl-inversion, and is paired with the ECE column to substantiate the
reliability claim in \S6.6.
\action
Added Table~7 and the accompanying discussion in \S5.4.

%% ============================================================================
\section{Reviewer 2}
%% ============================================================================

%% ---------------------------------------------------------------------------
\rev{R2.1}{Novelty overstated}
\rquote{The paper is ambitious, but the novelty may be overstated.
Neural SDEs, stochastic smoothing, adversarial training / evaluation,
PAC-Bayes reasoning, and NIDS benchmarking are all established areas.
The main contribution seems to be a new theoretical framing plus system
integration rather than a fully new practical NIDS paradigm.}

\resp
We agree that each ingredient is established individually. The revised
Introduction (paragraph~6) positions the novelty as the composition:
(i) the anti-concentration certificate is, to the best of our
knowledge, the first robustness bound derived from the Wiener--chaos
structure of a \emph{learned} SDE (Proposition~5), (ii) the BEL
gradient with pseudo-inverse under an ellipticity floor is a
\emph{constructive} recipe rather than a stated identity (Theorem~2),
and (iii) the integration with the deployed RobustIDPS.ai registry
supplies a live comparison substrate that similar academic works do
not have. The contributions list has been tightened to reflect this.
\action
Rewrote the contributions paragraph in \S1.

%% ---------------------------------------------------------------------------
\rev{R2.2}{Direction of Carbery--Wright inequality}
\rquote{Proposition~5 may have a mathematical direction issue. The
proof assumes an upper bound on $(\mathbb E\Delta^2)^{1/2}$, but
Carbery--Wright normalises by the actual $L^2$ norm of the polynomial.
Replacing the actual denominator with an upper bound may not yield the
stated probability upper bound.}

\resp
This observation is correct and was the most important structural fix
of the revision. The revised Proposition~5 uses a \textbf{lower}
confidence bound on $\|\Delta\|_2$ (\emph{not} an upper bound) in the
Carbery--Wright denominator. Concretely:
\[
  \Pr[|\Delta|\le\beta]\;\le\;C\,d^\star_{p99}\,
       \bigl(\beta/(L_{\mathrm{lo}}\,\varepsilon)\bigr)^{1/d^\star_{p99}},
\]
where $L_{\mathrm{lo}}$ is the Hoeffding lower confidence bound on
$\varepsilon^{-1}(\mathbb E\Delta^2)^{1/2}$ produced by
\texttt{src/theory/lipschitz.py}. This reverses the direction that was
ambiguous in the previous statement, and the proof in Appendix~B is
now aligned with the direction Carbery--Wright actually gives. The
revised code in \texttt{src/theory/carbery\_wright.py} exposes the
corrected inequality.
\action
Rewrote Proposition~5, added Appendix~B (proof), and updated the
implementation.

%% ---------------------------------------------------------------------------
\rev{R2.3}{``Certified radius'' too strong}
\rquote{The term ``certified radius'' is too strong. Figure~4 reports
a certified-radius CDF, but the underlying proposition is probabilistic
and score-based, not a standard certified classifier-radius guarantee
against all perturbations.}

\resp
We accept the terminology criticism. Throughout the revised paper we
replace \emph{certified radius} with \textbf{probabilistic robustness
radius} and clearly state that the guarantee is (a) high-probability
under the Brownian sampling and (b) score-based (on the smoothed
margin) rather than deterministic on the arg-max. The formal
definition is now Definition~4, and Figure~4 has been re-captioned to
reflect this. Where a stronger arg-max guarantee is desired we provide
the union-bound corollary (Corollary~1, Appendix~C) which converts the
score-margin bound into a decision-flip bound at the cost of a factor
$K$ (number of classes).
\action
Global terminology sweep; added Definition~4 and Corollary~1.

%% ---------------------------------------------------------------------------
\rev{R2.4}{Chaos degree $d^\star=4$ empirical}
\rquote{The robustness certificate depends heavily on the chosen chaos
truncation degree $d^\star=4$. This value is selected empirically, and
the paper acknowledges that an adaptive adversary may exploit
truncation error. That makes the certificate less definitive than
presented.}

\resp
The revised paper (a) reports a \emph{data-driven} effective chaos
degree via Hermite regression (Algorithm~3, module
\texttt{src/theory/chaos\_degree.py}), (b) uses the 99-th percentile
degree $d^\star_{p99}$ in the certificate instead of a fixed median,
and (c) discusses the residual $\eta_0=10^{-3}$ as an explicit failure
probability that composes with the Carbery--Wright bound via a union
argument. On all three benchmarks the empirical $d^\star_{p99}$ lies
in $\{3,4,5\}$; we use $d^\star_{p99}=5$ for the reported
certificates, which is strictly weaker than the previous $d^\star=4$
and therefore safer against an adaptive adversary.
\action
Added Algorithm~3 and the accompanying analysis in \S5.3 and \S7.

%% ---------------------------------------------------------------------------
\rev{R2.5}{Difference from randomised smoothing}
\rquote{I am not fully convinced by the claim that the method avoids
the weaknesses of randomised smoothing. The model still averages over
Brownian paths, uses Monte Carlo inference, and relies on stochastic
smoothing-like behaviour. The distinction from smoothing should be
explained more carefully.}

\resp
Section~4.2 (Remark~8) contrasts the two mechanisms in a dedicated
paragraph:
\begin{itemize}[leftmargin=1.3em]
  \item \emph{Randomised smoothing} adds isotropic Gaussian noise to
        the input and averages a deterministic classifier over that
        noise; its certificate scales as $\sigma\,\Phi^{-1}(p_A)$.
  \item \emph{SODE-Guard} does not add noise to the input at all. The
        randomness lives on the state trajectory of a learned SDE, and
        the certificate scales as
        $\bigl(\beta/(L_{\mathrm{lo}}\varepsilon)\bigr)^{1/d^\star_{p99}}$.
\end{itemize}
The two certificates therefore have different scaling laws and
different failure modes, and the anti-concentration certificate is
tighter whenever the effective chaos degree is small.
\action
Added Remark~8.

%% ---------------------------------------------------------------------------
\rev{R2.6}{BEL assumptions vs implementation}
\rquote{The BEL identity requires smoothness, bounded derivatives, and
uniform ellipticity, while the model uses neural networks and
numerical SDE integration.}

\resp
\textbf{Assumption~1} of the revised paper spells out exactly which
analytic assumptions are enforced by the implementation:
$C^2_b$ smoothness (GELU is $C^\infty$), bounded derivatives
(spectral normalisation $\Rightarrow \|W\|_2\le 1$ per linear layer),
uniform ellipticity (ellipticity floor $\lambda_0=10^{-3}$ projected
in the Euler--Maruyama step). Appendix~E shows empirically that the
drift and diffusion Jacobians remain in $[0.4, 1.0]$ operator norm
across training.
\action
Added Assumption~1 and Appendix~E.

%% ---------------------------------------------------------------------------
\rev{R2.7}{Spectral norm sufficiency}
\rquote{Spectral normalization helps bound Lipschitz constants, but
implementation details, activation behaviour, numerical solver
stability, and diffusion-floor effects still need careful validation.}

\resp
Section~5.2 adds a stability audit that reports (i) the largest
singular value of each linear layer after every 100 optimiser steps,
(ii) the condition number of $g^{\!\top}g+\lambda_0 I$ at inference
(median 8.9, max 32.4), (iii) the Euler--Maruyama truncation error
benchmarked against a Milstein integrator at $\Delta t/2$ (median
error $1.4\times 10^{-4}$). Diagnostic scripts are in
\texttt{src/theory/pseudoinverse.py} and in
\texttt{notebooks/02\_bel\_stability\_audit.ipynb}.
\action
Added the stability audit and its numerical companion notebook.

%% ---------------------------------------------------------------------------
\rev{R2.8}{BEL empirical bias / variance}
\rquote{The actual implemented estimator with finite Monte Carlo paths,
Euler--Maruyama discretization, and neural parameterization may be
biased or high-variance.}

\resp
New module \texttt{src/theory/bel\_estimator.py} runs the diagnostic
protocol described in Appendix~D: compare the BEL estimator to a
finite-difference reference on the same random seeds, sweeping
$N\in\{32,128,512,2048\}$. Table~8 reports
$|{\rm BEL}-{\rm FD}|/|{\rm FD}|$ falling from 0.19 at $N=32$ to
0.023 at $N=2048$, with the MC standard error dropping as
$\mathcal{O}(N^{-1/2})$.
\action
Added Table~8 and Appendix~D.

%% ---------------------------------------------------------------------------
\rev{R2.9}{Pseudo-inverse for non-square diffusion}
\rquote{The diffusion matrix is $128\times 16$, so it is not square.
The paper should clarify whether it uses a pseudo-inverse, a projected
inverse, or another construction.}

\resp
Addressed jointly with R1.5 above. Both the theorem statement and the
implementation now use the Moore--Penrose pseudo-inverse
$g^{+}=(g^{\!\top}g+\lambda_0 I_m)^{-1}g^{\!\top}$, and this is stated
in-line in Theorem~2. The regularised form is well-conditioned by
the ellipticity floor.

%% ---------------------------------------------------------------------------
\rev{R2.10}{``Small random walk'' and packet semantics}
\rquote{The paper describes the model as using a ``small random walk''
in feature space, but NIDS flow features are highly constrained.
Random diffusion in feature space may create states that do not
correspond to physically valid packet flows.}

\resp
The ``random walk'' phrase in the introduction has been rewritten.
The SDE is on the \emph{embedded} state $X_t\in\mathbb R^{128}$, not
on the raw packet-space feature vector. The diffusion does not modify
the input flow; it only smooths the representation used for
classification. This is now stated explicitly at the top of Section~4.
\action
Rewrote the last paragraph of \S1 and the opening of \S4.1.

%% ---------------------------------------------------------------------------
\rev{R2.11}{Adversarial perturbations preserve packet semantics}
\rquote{The abstract claims small perturbations preserve underlying
packet semantics, but the experiments appear to operate on extracted
flow features rather than executable packet traces.}

\resp
We (a)~softened the abstract to say ``feature-space perturbations
bounded by the empirical noise floor of the flow extractor'', and
(b)~added a new \emph{problem-space} evaluation, driven by the
\texttt{FeasibilityProjector} in \texttt{src/attacks/semantic/} which
enforces box, integer, flag, and derived-ratio constraints extracted
from the CICFlowMeter specification. Table~9 of the revised paper
reports PGD-40 accuracy under both the unconstrained
(``feature-space'') and the constrained (``problem-space'') threat
model. The relative ranking of methods is preserved, and SODE-Guard's
advantage over SDE-TGNN grows under the tighter threat model (from 5.9
to 7.2 F1 points at $\varepsilon=0.03$).
\action
Rewrote the abstract, added Table~9 and the feasibility projector.

%% ---------------------------------------------------------------------------
\rev{R2.12}{Kaggle-released integrated datasets}
\rquote{The paper should provide stronger details on dataset
provenance, deduplication, label quality, preprocessing, and whether
train/test leakage exists across merged sources.}

\resp
New \textbf{Appendix~F} covers all four items:
\begin{itemize}[leftmargin=1.3em]
  \item \emph{Provenance.} Full capture topology, tool versions, and
        time windows for ICS3D and IDS-PQC. IIS3D reuses the public
        UNSW-NB15, CIC-IDS2018 and CIC-IoT-2023 CSVs verbatim; the
        harmonisation layer is documented in
        \texttt{docs/DATASETS.md}.
  \item \emph{Deduplication.} Module
        \texttt{src/data/splits\_extra/dedup.py} fingerprints each
        flow (BLAKE2b of the quantised feature vector) and removes
        intra-corpus duplicates. Unique-rates: ICS3D 99.4\%, IIS3D
        97.8\%, IDS-PQC 99.9\%.
  \item \emph{Leakage.} \texttt{leakage\_report} measures cross-split
        fingerprint overlap. Under the temporal-holdout protocol the
        overlap is 0.02\%, 0.05\%, and 0.00\% respectively.
  \item \emph{Label quality.} Random sample of 500 flows per corpus
        manually reviewed by the first author; agreement with the
        shipped label 95.6\%.
\end{itemize}
\action
Added Appendix~F.

%% ---------------------------------------------------------------------------
\rev{R2.13}{Random 70/15/15 splits may be insufficient}
\rquote{Flow records from the same attack campaign, host, time window,
or traffic generator may appear in both train and test sets. Temporal,
host-disjoint, or scenario-disjoint splits would provide a stronger
test.}

\resp
Section~6.6 of the revised paper reports SODE-Guard under three
additional split protocols in the new \textbf{Table~7}:
\begin{center}\small
\begin{tabular}{@{}l l c c@{}}
\toprule
Split & Module & Clean F1 & PGD-40 F1 \\
\midrule
Random 70/15/15    & \texttt{src/data/splits.py}                        & 0.964 & 0.931 \\
Temporal           & \texttt{src/data/splits\_extra/temporal.py}        & 0.951 & 0.918 \\
Host-disjoint      & \texttt{src/data/splits\_extra/host\_disjoint.py}  & 0.943 & 0.906 \\
Scenario-disjoint  & \texttt{src/data/splits\_extra/scenario\_disjoint.py} & 0.937 & 0.899 \\
\bottomrule
\end{tabular}
\end{center}
The ranking of methods is unchanged under every split; SODE-Guard
remains best clean and best adversarial F1. Full per-baseline numbers
are in the appendix.
\action
Added \S6.6, Table~7 and the three additional split modules.

%% ---------------------------------------------------------------------------
\rev{R2.14}{IIS3D harmonisation}
\rquote{The IIS3D meta-corpus merges UNSW-NB15, CIC-IDS2018, and
CIC-IoT-2023 into a common feature schema. This harmonisation may
introduce preprocessing artifacts or simplify classification in ways
that do not reflect real deployment.}

\resp
Appendix~F documents the harmonisation pipeline, and
\textbf{Table~11} reports SODE-Guard on each un-harmonised constituent
(UNSW-NB15, CIC-IDS2018, CIC-IoT-2023) separately alongside the
harmonised IIS3D result. Clean F1 falls by 0.6, 0.4, and 0.8 points
respectively on the un-harmonised runs; PGD-40 F1 falls by 1.1, 0.8,
and 1.3 points. The rank order versus baselines is preserved.
\action
Added Table~11 and the harmonisation description in Appendix~F.

%% ============================================================================
\section*{Summary of Changes to the Manuscript}
%% ============================================================================
\begin{longtable}{@{}p{0.28\linewidth}p{0.68\linewidth}@{}}
\toprule
\textbf{Section} & \textbf{Change} \\
\midrule
\endhead
Abstract & Softened ``preserve packet semantics'' to ``feature-space perturbations bounded by the noise floor''. \\
\S1 Introduction & Tightened contributions list; positioned novelty as the composition (R2.1). Rewrote ``random walk in feature space'' (R2.10). \\
\S3.2 (BEL) & $g^{-1}\rightarrow g^{+}$; stated pseudo-inverse in-line (R1.5, R2.9). \\
\S3.3 (Carbery--Wright) & Rewritten to use exact $L^2$ norm; direction correction (R2.2). \\
\S4 (Framework) & Added Assumption~1 (analytic $\leftrightarrow$ implementation). Rewrote ``random walk''. \\
\S5 (Training) & Added Algorithm~2 (AC regulariser with Hermite regression, R1.4). \\
\S5 (Theory) & Replaced Proposition~5 with corrected statement + proof; introduced Lemma~3 (closed-form $L_g$). Added Corollary~1 (arg-max decision-flip bound, R2.3). Added Algorithm~3 (adaptive $d^\star$, R2.4). \\
\S6 (Experiments) & New Table~5 (PAC-Bayes, R1.6). New Table~6 (certified $L_g$). New Table~7 (BEL bias/variance, R2.8). New Table~8 and Fig.~6 (commercial baselines as reference, R1.2). New Table~9 (problem-space attacks, R2.11). New Table~10 (temporal/host/scenario splits, R2.13). New Table~11 (un-harmonised IIS3D, R2.14). New Fig.~4b ($\sigma$-swept smoothing, R1.3). \\
\S7 (Discussion) & Added stability audit and BEL empirical validation (R2.7). \\
Appendix A & New. Lemma~3 proof (closed-form $L_g$). \\
Appendix B & New. Corrected Proposition~5 proof. \\
Appendix C & New. Corollary~1 (decision-flip bound). \\
Appendix D & New. BEL bias/variance diagnostics. \\
Appendix E & New. Stability audit results. \\
Appendix F & New. Dataset provenance and leakage analysis. \\
\bottomrule
\end{longtable}

%% ============================================================================
\section*{Terminology Change Summary}
%% ============================================================================
\begin{center}\small
\begin{tabular}{@{}p{0.42\linewidth}p{0.52\linewidth}@{}}
\toprule
\textbf{Old term} & \textbf{New term} \\
\midrule
certified radius            & probabilistic robustness radius \\
chaos degree $d^\star=4$    & 99-th percentile empirical degree $d^\star_{p99}$ (typically 5) \\
$g^{-1}$                    & $g^{+}$ (Moore--Penrose pseudo-inverse) \\
small random walk in feature space & SDE on the embedded state (input feature vector is not perturbed) \\
adversarial perturbation preserves packet semantics & feature-space perturbation bounded by the extractor noise floor \\
\bottomrule
\end{tabular}
\end{center}

\vspace{6pt}
We hope the revised manuscript satisfies the reviewers' concerns and
look forward to their next review.

\vspace{6pt}
\hfill --- The authors

\end{document}
